\begin{abstract}
Cloud-native NoSQL databases such as Amazon DynamoDB have emerged as a foundational component of serverless architectures, yet the performance–cost trade-offs introduced by partition-key design and capacity mode remain empirically underexplored. This study presents a factorial experimental evaluation that simultaneously varies partition-key design (simple, composite, and write-sharded keys) and capacity mode (on-demand and provisioned with auto-scaling) across four serverless workload profiles driven by AWS Lambda. Building on the latency and throughput measures established by \citet{pantelic2026benchmarking} for self-hosted NoSQL persistence under microservice workloads, this research extends the evaluation to metered infrastructure by quantifying throttling rates, consumed read and write capacity units (RCU/WCU), and absolute cost per 10,000 operations—dependent variables absent from self-hosted baselines. The artefact comprises Infrastructure-as-Code (Terraform), Lambda-based workload generators with Zipfian access distributions, CloudWatch metrics collectors, and statistical analysis scripts implementing two-way ANOVA with interaction effects. Experimental results, captured through controlled moto-based simulation and designed for validation on live AWS infrastructure, demonstrate the performance and cost implications of configuration choices under sustained and burst load conditions. Findings provide empirical guidance for cloud architects selecting DynamoDB table configurations in latency-sensitive, cost-constrained serverless deployments.
\end{abstract}
